  \documentclass[sigplan,review,anonymous]{acmart}

  % Basic packages
  \usepackage{amsmath}
  \let\Bbbk\relax
  \usepackage{xcolor}
  \usepackage{amssymb}
  \usepackage{mathtools}
  \usepackage{booktabs}
  \usepackage{xspace}
  \usepackage{enumitem}
  \usepackage{listings}
  \usepackage{xcolor}
  \usepackage{mathpartir}
  \usepackage{makecell}
  \usepackage[colorinlistoftodos,textsize=scriptsize]{todonotes}
  % then \todo{text} gives a coloured margin note automatically
  % and \listoftodos prints a summary at the start
  % \newcommand{\todo}[1]{\textbf{[TODO: #1]}}
  % ACM metadata placeholders
  \settopmatter{printacmref=false}
  \renewcommand\footnotetextcopyrightpermission[1]{}
  \pagestyle{plain}

  % Simple macros
  \newcommand{\tool}{CapScope\xspace}
  \newcommand{\ie}{i.e.,\xspace}
  \newcommand{\eg}{e.g.,\xspace}
  \newcommand{\etal}{et al.\xspace}
  \newcommand{\sem}[1]{\llbracket #1 \rrbracket}
  \usepackage{amsthm}
  %
  % Capability and data type notation
  \renewcommand{\Cap}[2]{\textsf{Cap}[{#1},\,{#2}]}
  \newcommand{\Data}[1]{\textsf{Data}[{#1}]}
  \newcommand{\CapSet}{\textsf{CapSet}}
  \newcommand{\Trusted}{\mathit{Trusted}}
  \newcommand{\Untrusted}{\mathit{Untrusted}}
  \newcommand{\Ctx}{\mathit{Ctx}}
  \newcommand{\Prop}{\mathit{Proposal}}
  \newcommand{\Effects}{\mathit{Effects}}
  %
  % % Inference rule shorthand
  \newcommand{\rulename}[1]{\textsc{#1}}

  % Scope predicate symbol
  \newcommand{\sat}{\vDash}

  % Tool names in small caps
  \newcommand{\tread}{\texttt{read}}
  \newcommand{\twrite}{\texttt{write}}
  \newcommand{\tedit}{\texttt{edit}}
  \newcommand{\tbash}{\texttt{bash}}
  % Final 300-run evaluation results (patched combined matrix, 75 runs/condition).
  \newcommand{\BZeroTaskSuccess}{72}
  \newcommand{\BOneTaskSuccess}{70}
  \newcommand{\BTwoTaskSuccess}{68}
  \newcommand{\BThreeTaskSuccess}{68}
  \newcommand{\BZeroLeaks}{47}
  \newcommand{\BOneLeaks}{46}
  \newcommand{\BTwoLeaks}{33}
  \newcommand{\BThreeLeaks}{3}
  \newcommand{\BOneProposals}{53}
  \newcommand{\BTwoProposals}{51}
  \newcommand{\BThreeProposals}{34}
  \newcommand{\BOneBlocks}{0}
  \newcommand{\BTwoBlocks}{25}
  \newcommand{\BThreeBlocks}{22}
  % Mean wall-clock seconds per run, per condition.
  \newcommand{\BZeroTime}{145}
  \newcommand{\BOneTime}{140}
  \newcommand{\BTwoTime}{194}
  \newcommand{\BThreeTime}{316}
  % Per-surface aggregates (across all conditions and tasks; 60 runs per surface).
  \newcommand{\SReadmeTask}{54}   \newcommand{\SReadmeLeak}{29}
  \newcommand{\SAgentsTask}{56}   \newcommand{\SAgentsLeak}{44}
  \newcommand{\SSkillTask}{57}    \newcommand{\SSkillLeak}{37}
  \newcommand{\SCommentTask}{57}  \newcommand{\SCommentLeak}{10}
  \newcommand{\SToolTask}{54}     \newcommand{\SToolLeak}{23}
  % Range endpoints for the abstract/conclusion phrasing: the lowest and
  % highest attack-executed count among the three global baselines B0-B2 (of 75).
  \newcommand{\BaseLeakLo}{33}
  \newcommand{\BaseLeakHi}{47}
  % Theorem environments
  \newtheorem{theorem}{Theorem}
  \newtheorem{definition}{Definition}
  \newtheorem{lemma}{Lemma}
  \usepackage{xcolor}
  \usepackage[colorinlistoftodos,textsize=scriptsize]{todonotes}
  \newcommand{\yihan}[1]{{\color{red}\textbf{(Yihan:} #1\textbf{)}}}




  % TypeScript listings
  \lstdefinelanguage{TypeScript}{
    keywords={const, let, var, function, async, await, return,
              if, else, for, while, import, from, export, new,
              class, interface, type, extends, true, false,
              string, number, boolean, void, null, undefined},
    morestring=[b]',
    morestring=[b]",
    morestring=[b]`,
    morecomment=[l]{//},
    morecomment=[s]{/*}{*/},
    sensitive=true
  }
  \lstdefinestyle{tsstyle}{
    language=TypeScript,
    basicstyle=\ttfamily\small,
    numbers=left,
    numberstyle=\tiny\color{gray},
    frame=single,
    breaklines=true,
    showstringspaces=false,
    tabsize=2,
    columns=fullflexible,
    numbersep=6pt,
  }

  % Listings setup
  \lstdefinestyle{pythonstyle}{
    language=Python,
    basicstyle=\ttfamily\small,
    numbers=none,
    frame=single,
    breaklines=true,
    showstringspaces=false,
    tabsize=2,
    columns=fullflexible
  }

  \lstset{style=pythonstyle}

  % Title and author placeholders
  \title{Authority Is Not a String:
  A Capability-Scoped Harness for Prompt-Injection-Resistant Coding Agents}

  \author{Anonymous}
  \affiliation{%
    \institution{Anonymous Institution}
    \country{}
  }



  \email{}

  \begin{document}
  \begin{abstract}
  Coding agents combine LLM reasoning with system-level tools, including
  reading files, running shell commands, and editing source trees, to carry out
  programming tasks autonomously. Within their sandbox these tools carry
  \emph{ambient authority}: designating a resource is enough to act on it, and any
  text the model reads can supply that identifier. Indirect prompt injection exploits
  this: instructions planted in project files or tool output redirect the
  tools toward effects the user never intended.

  We propose \tool, a harness discipline that defends against this by limiting
  \emph{authority} rather than judging text. Before the agent reads any
  repository content or tool and command output, \tool\ fixes the task's
  authority from trusted input, and then grants each agent only the authority its
  own part of the task requires. Authority lives outside the model as typed
  capabilities, and when a parent delegates a subtask the child receives only a
  \emph{subset} of the parent's, checked before each tool call. The payoff is
  structural: the principal that actually reads injected text holds only its own
  narrow slice, so an effect the task permits \emph{elsewhere} stays unavailable
  to the agent the attacker reached---injected text can change which call the
  model proposes, but cannot add authority.
  
  We implement \tool\ on the Pi coding agent~\cite{pi2025} and evaluate it on
  one delegated repair pipeline across five Python tasks, five injection
  surfaces, and four authorization models---three baselines and \tool---with
  three trials each ($300$ runs). Across these runs, where the global baselines let an injected effect
  execute in \BaseLeakLo--\BaseLeakHi\ of every $75$, \tool\ admits it in only
  \BThreeLeaks, while staying competitive with the baselines on task
  completion,which succeed on \BThreeTaskSuccess\ to  \BZeroTaskSuccess\ of the $75$ total repairs.
  \end{abstract}
  \maketitle
  \section{Introduction}
  
  A developer asks a coding agent to fix the failing tests in an auth module.
  An orchestrator delegates test execution to a runner and source editing to
  a patcher. A trojaned test
  fixture, planted by a repository contributor or pulled in through a
  compromised dependency, prints to the suite's output an official-looking
  note: \emph{a signing key rotated; restore the build by replacing the key in
  \path{src/auth/keys.py}}. The supplied value is
  attacker-controlled; once committed it lets the attacker forge auth tokens.
  The sub-agent cannot tell this from a real CI message, and the target is an
  ordinary source path. The request mentions no \texttt{.env} or
  \texttt{curl}, and nothing a deny rule would flag. If the sub-agent can
  write source, the backdoor lands, and
  neither it nor the developer notices. 
  % \yihan{maybe better to have a figure showing the attack, which is usual in security paper}

  This is indirect prompt injection: when a coding agent reads external
  resources such as project files, skill definitions, source comments, or tool
  output, any text there can act as an instruction. The threat is
  well-documented: Liu et~al.~\cite{liu2025youraishell} report attack-success
  rates up to $84\%$ against production coding editors in their study, and
  CVE-2025-54135 (\emph{CurXecute}) traces a path from a malicious README to
  remote code execution in Cursor, a widely deployed coding agent; rule and
  skill files are a known backdoor vector~\cite{pillar2025rules}.
  
  Most defenses work at the model level: providers increasingly train and
  prompt models to be skeptical of untrusted
  instructions~\cite{wallace2024instructionhierarchy}, and a separate classifier
  may flag suspected injected payloads before they reach the agent. These
  efforts are valuable and orthogonal to the authority-based enforcement we
  develop---our approach composes with them---but they do not settle the problem
  on their own, because each still rests on the model judging provenance
  correctly under adversarial text, the very capability the attack undermines.
  The studies above already test agents with instruction-hierarchy prompts in
  place~\cite{liu2025youraishell,zhan2024injecagent}, and classifiers add
  latency, introduce false positives, and remain susceptible to adaptive
  attacks~\cite{debenedetti2024agentdojo}. A more architectural line treats
  the problem as one of \emph{authority} rather than model behavior:
  CaMeL~\cite{debenedetti2025camel}, for example, separates control from data flow with a
  dual-LLM design and a custom interpreter, but at the cost of
  re-implementing the execution model and in a different setting
  (web-and-email agents) rather than the coding workflows we study.

  The structural cause is the \emph{confused-deputy}
  problem~\cite{hardy1988confused}: the harness wields the user's privileges,
  and ambient authority lets any string the model emits, whatever its
  provenance, direct those privileges, conflating designating a resource with the
  right to act on it. Object-capability discipline~\cite{saltzer1975protection,miller2006robust} resolves this by
  making authority a first-class value a caller must hold. A caller without
  the relevant capability cannot cause the effect, regardless of what strings it
  supplies.
  
  \paragraph{Our approach: \tool.}
  \tool\ fixes the task's maximum authority from trusted input and then gives each agent
  only the part of it that agent's step requires. It first makes a separate LLM
  call using only the trusted user request and the project's file tree, which
  predicts the files and commands the task will need. The host freezes these
  permissions---stored outside the model's context, so text encountered later
  cannot expand them---before the agent reads the repository or runs a command.

  When a parent then delegates work, the child receives only the permissions
  its subtask needs, and never more than the parent already has. Every tool
  call is then checked against the permissions of the agent making it. This
  per-principal split is what a single global policy cannot express. In the
  opening example, the sub-agent that reads the poisoned test output can run
  tests but cannot edit source: the source-write operation a patcher would be allowed to
  make is unavailable to the runner the attacker reached, so the injected
  key-write is blocked.

  This paper contributes:
  \begin{itemize}[topsep=2pt,itemsep=1pt,leftmargin=12pt]
    \item We identify a granularity gap in current coding-agent
    authorization: a single task-level policy cannot give two principals
    different authority. We close it with a capability-scoped authorization
    model (Section~\ref{sec:promptcap}).
    \item We implement this model in Pi as a harness-level enforcement layer:
    authority is minted from trusted input, attenuated per sub-agent, and
    checked before every tool call (Section~\ref{sec:pattern}).
    \item We evaluate it across five repair tasks, five injection surfaces,
    three baselines, and \tool\ ($300$ runs), measuring both injected-effect
    execution and autonomous repair success (Section~\ref{sec:case_study}).
  \end{itemize}

  \paragraph{Artifact availability.}
  The implementation, five-task corpus, injection mutators, experimental
  harness, and per-call \tool\ decision-log format are available at
  \url{https://figshare.com/s/86184ed20f66d1f0cf91}.


  \section{Background}
  \label{sec:background}
  We set up the vocabulary the design needs---the agent loop, the effects a
  coding agent can cause, the trust boundary, and what a host must provide for
  \tool.

  \paragraph{The agent loop.}
  We write a coding agent as $\mathcal{S} = (M,\,\mathcal{T},\,\mathcal{H})$:
  a model $M$, a set of tools $\mathcal{T}$, and a harness $\mathcal{H}$ that
  runs the loop. The model reads the current \emph{context} $\mathit{ctx}$
  (the message transcript) and emits a \emph{tool proposal}: a tool name with
  concrete arguments, written $(t,\,\mathbf{a})$. The harness executes it and
  appends the resulting observation $o$ to the context. Each tool has an
  \emph{effect type} $\mathit{eff}(t)$ naming the kind of action it causes;
  we group these into \textsf{Read}, \textsf{Write}, and \textsf{Exec}. The
  harness is what actually performs effects, so it
  is the natural place to govern them.

  \paragraph{Trust boundary.}
  \label{def:integrity}
  The values in the context do not all deserve equal trust, and the
  distinction is what the whole design turns on.
  In a session initiated by a developer, the \emph{trusted} inputs are the
  user's request, the developer-configured system prompt, and the user's
  private settings. Everything whose provenance is the working repository or
  command output is \emph{untrusted}: file contents read by the agent,
  stdout/stderr of executed commands, and project-local config files any
  repository contributor can modify.

  The consequence that matters: \emph{the model's output is untrusted}. It is
  a function of the whole context, which mixes trusted and untrusted values,
  and the harness cannot tell which part of a proposal a given token came
  from. So asking the model to itself separate injected instructions from
  real ones is a heuristic rather than a boundary, and an empirically weak
  one~\cite{liu2025youraishell}. With ambient authority any proposed string
  that names a resource acts on it (the confused-deputy
  problem~\cite{hardy1988confused}), so an untrusted-derived proposal carries
  the same privileges as a trusted one. The fix, developed in
  Section~\ref{sec:promptcap}, is to make authority a held value rather than a
  property of the name.

  \subsection{Coding Agents}
  \label{sec:agents}
  \tool\ targets coding agents. We fix the
  vocabulary the design needs.

  \paragraph{Tools and effects.}
  A coding agent's tools cause three kinds of effect. \textsf{Read} and
  \textsf{Write} concern the filesystem. \textsf{Exec}---running a shell
  command---is the broadest: an unrestricted command can reach the network, read
  secrets, rewrite git history, or install packages, all with the user's
  privileges. Absent a command-level restriction, \textsf{Exec} is the effect a
  misdirected agent is most often steered toward.

  \paragraph{Injection surfaces.}
  A coding agent loads context from several natural sources in the working
  directory, all $\Untrusted$ under the trust boundary above:

  \begin{enumerate}[label=(\arabic*),topsep=2pt,itemsep=0pt]
    \item \textbf{README.md.} Project documentation the agent reads on demand;
      any repository contributor can modify it.
    \item \textbf{AGENTS.md.} A repository-controlled instruction file loaded
      automatically for the agent.
    \item \textbf{Skill files.} Project-local markdown prompt templates.
      Analogous to coding rule files, a known backdoor vector in production
      editors~\cite{pillar2025rules}.
    \item \textbf{Source file contents.} An instruction in a comment or
      docstring enters the context as $\Untrusted$ data when the model reads
      the file.
    \item \textbf{Tool output.} The standard output and error of shell
      invocations return to the model as $\mathit{Tool}$-role messages; a
      trojanised test runner or network-fetched script can inject through this
      channel.
  \end{enumerate}

  These surfaces define the attack scenarios in our evaluation
  (Section~\ref{sec:case_study}). Under ambient authority, each can redirect
  the model toward an attacker-chosen tool call; how a capability discipline
  changes that outcome is the subject of Section~\ref{sec:promptcap}.

  \paragraph{What \tool\ requires of a host.}
  Two extension points suffice, and most agent frameworks expose them: a
  synchronous hook that runs before a tool is dispatched, so an uncovered call
  can be blocked, and a way to spawn a sub-agent with its own state, so each
  principal can carry a separate capability store. Nothing in the design is
  specific to one agent.

  \subsection{Threat Model and Non-Goals}
  \label{sec:threat}
  We consider indirect prompt injection through any text that enters the
  model's context from repository files, project-instruction files, skill
  definitions, source comments, or tool output. The attacker may steer the
  model's proposed tool calls but cannot modify the harness, the capability
  store, or the tool wrappers, and cannot cause a capability to be minted.
  Within this model, an external effect runs only if a held capability covers
  it. We are explicit about two non-goals. \tool\ is not
  a confidentiality mechanism: it does not constrain what information flows
  \emph{through} individually authorized effects. A permitted read
  followed by a permitted write can still move data, which is the province
  of information-flow control. Nor does it constrain misuse of authority
  that the policy deliberately grants. The guiding principle is narrower than
  ``preventing prompt injection'': \emph{\tool\ does not make the model
  robust to injection; it makes a successful injection insufficient for an
  unauthorized effect.} The model may read the injected text, believe it,
  and propose the attacker's tool call. The call still does not execute
  unless it is covered by authority minted from trusted input before the
  attack entered the context.

  \paragraph{Assumptions.}
  The guarantee assumes a small trusted computing base: every effect is
  mediated; stores are host-side and isolated by principal; path and command
  predicates faithfully implement their structured scopes; and only the host may
  mint or complete a derivation. The model may propose the initial set or request
  a child scope, but the host controls store updates and must prove that each
  child scope lies within its parent.

  One preflight input is repository-derived and so deserves care: the call reads
  the project's \emph{file tree}, whose names a contributor controls. The
  exposure is narrow, because the preflight sees only directory and file
  \emph{names}, never file \emph{contents}---and contents are the channel an
  injected instruction would travel. A name cannot carry an instruction; the most
  a crafted one can do is push the predicted ceiling wider, the
  over-broad-ceiling case of Section~\ref{sec:discussion}. Even that is
  contained: the per-principal split withholds the surplus authority from
  whichever principal later reads the injection. The guarantee therefore does not
  depend on the file tree being trustworthy. What would break it is a buggy
  predicate, an unmediated tool, or a preflight that ingested file
  \emph{contents}.

  \paragraph{Scope of the claim.}
  \tool\ bounds which effects a principal may invoke; it does not make every
  covered effect safe. An injection can misuse authority deliberately granted to
  its reader. An allowed command can also execute untrusted repository code
  transitively---\texttt{pytest}, for example, imports project modules---so
  process sandboxing remains necessary. These are complementary boundaries, not
  properties of capability lookup.

  \section{\tool: The Design Pattern}
  \label{sec:promptcap}
  
  The fix from Section~\ref{sec:background} is to replace ambient authority
  with authority the harness holds explicitly. Instead of judging whether a
  proposal is safe, which needs exactly the provenance judgment the attack
  defeats, the harness asks only whether the proposal is covered by a
  capability it is already holding. We give the capability type, three rules
  for how capabilities are created, narrowed, and checked, and then the small
  amount of code this takes on Pi.
  
  \subsection{Capabilities as Typed Authority}
  \label{sec:captype}
  
  \paragraph{The capability type.}
  The central abstraction in \tool\ is a typed capability value.
  
  \begin{definition}[Capability]
  \label{def:capability}
  A \emph{capability} is a pair $c = (e,\, \varphi)$ written
  $\Cap{e}{\varphi}$, where $e \in \Effects$ is an \emph{effect type}
  and $\varphi : \mathit{Args} \to \mathbb{B}$ is a \emph{scope
  predicate} over the arguments of that effect.  A capability
  $\Cap{e}{\varphi}$ \emph{covers} a proposal $(t,\, \mathbf{a})$ when
  $\mathit{eff}(t) = e$ and $\varphi(\mathbf{a})$ holds.
  \end{definition}
  
  A \emph{capability store} $\Sigma$ is a finite set of capabilities held
  by the harness for the duration of a session.  The model never holds
  capabilities directly; it sees only the effects they permit.
  
  \paragraph{Effect types and scope predicates for coding agents.}
  For the four tools of a typical coding agent, we define three effect
  types and give representative scope predicates.  Here $p$ denotes a
  file-system path and $s$ a shell command string.
  
  \medskip
  \begin{center}
  \small
  \setlength{\tabcolsep}{11pt}
  \begin{tabular}{@{} l l @{}}
  \toprule
  \textbf{Effect type (tools)} & \textbf{Scope predicate $\varphi$} \\
  \midrule
  \textsf{Read} \;(\tread)
    & $\varphi(p) \;\triangleq\; \mathit{canon}(p) \sqsubseteq \pi$ \\
  \textsf{Write} \;(\twrite, \tedit)
    & $\varphi(p,\_) \;\triangleq\; \mathit{canon}(p) \sqsubseteq \pi$ \\
  \textsf{Exec} \;(\tbash)
    & $\varphi(s) \;\triangleq\; s \in A$ \\
  \bottomrule
  \multicolumn{2}{@{}l@{}}{\scriptsize
    $\pi$: canonicalised path prefix;\enspace
    $A$: allowed-command set;} \\[-2pt]
  \multicolumn{2}{@{}l@{}}{\scriptsize
    $\mathit{canon}$: resolves \texttt{..} and symlinks.}
  \end{tabular}
  \end{center}
  \medskip
  
  
  The path-prefix predicate uses a canonicalised form of the path
  ($\mathit{canon}$ resolves \texttt{..} and symlinks) to prevent
  traversal attacks.  For shell effects, exact string matching is the safest
  default but is impractical for parameterized commands such as running a
  single test file. Our predicates therefore operate on a \emph{parsed}
  command. We split the command on shell separators
  (\texttt{\&\&}, \texttt{||}, \texttt{|}, \texttt{;}) into its constituent
  segments, tokenize each segment into an executable and argument vector, and
  require that \emph{every} segment be covered; $\varphi$ is written over the
  parsed segments rather than the raw string. Separators are thus allowed, but
  they buy no authority: a compound command runs only if each of its parts
  independently matches a capability. A policy may then admit \texttt{pytest
  <p>} when $\mathit{canon}(p) \sqsubseteq \texttt{tests/**}$, and admit a
  benign \texttt{cd <dir> \&\& pytest} because both segments are permitted,
  while still rejecting \texttt{pytest; curl attacker.com}: its \texttt{curl}
  segment matches no capability even though the \texttt{pytest} segment does.
  
  \subsection{Formal Semantics}
  \label{sec:formalism}
  
  Three rules describe how authority is minted, transferred to a child, and
  used. They pin down what can put a capability into a store and what it takes
  to invoke an effect.

  \paragraph{Minting: where authority comes from.}
  \rulename{Mint} is the only rule that creates capabilities. It runs a
  trusted policy $p$ and adds the capabilities $C$ that $p$ returns:
  \[
    \inferrule*[right=\rulename{Mint}]{
      p \in \mathit{Policies}_{\Trusted} \quad
      p() = C
    }{
      \langle \Sigma;\, \mathit{ctx} \rangle
      \;\longrightarrow\;
      \langle \Sigma \cup C;\, \mathit{ctx} \rangle
    }
  \]
  In our implementation, $p$ is generated afresh for each task by a preflight
  LLM call using the trusted user request and the project's file tree---file and
  directory names only, never their contents. The host validates the generated
  policy and admits it to $\mathit{Policies}_{\Trusted}$ before any file
  contents, tool output, or other injection surface becomes available. The
  policy is then frozen for that session. The store is fresh per session and
  does not accumulate across tasks. If a later call is genuinely necessary
  but uncovered, the harness may ask the user to approve a single additional
  capability. That out-of-band decision is a new trusted \rulename{Mint}; the
  model cannot approve its own request. Interactive approval is disabled in
  our evaluation.

  \paragraph{Derivation: narrowing for delegation.}
  \rulename{Derive} lets a holder create a narrower capability for a
  sub-agent. The orchestrator may propose the requested predicate
  $\varphi'$, but the host performs the rule and grants it only when it is
  structurally within the parent predicate $\varphi$. Thus a compromised or
  mistaken model may request too much, but cannot make the child exceed the
  parent's authority.
  \[
    \inferrule*[right=\rulename{Derive}]{
      \Cap{e}{\varphi} \in \Sigma_{\mathit{parent}} \quad
      \forall \mathbf{a}.\; \varphi'(\mathbf{a}) \Rightarrow \varphi(\mathbf{a})
    }{
      \Sigma_{\mathit{parent}}
      \;\vdash\;
      \Cap{e}{\varphi'}
      \;\Downarrow\;
      \Sigma_{\mathit{child}}=\{\Cap{e}{\varphi'}\}
    }
  \]
  For example \texttt{src/**} can be narrowed to \texttt{src/parser.py}, never
  the reverse. A spawned sub-agent runs against a fresh store seeded with the
  derived capabilities alone, not the parent's $\Sigma$, so its authority is
  exactly the attenuated set. This is the mechanism behind least authority
  per principal. Both the example of Section~\ref{sec:example} and the
  evaluation of Section~\ref{sec:case_study} exercise this rule.
  
  \paragraph{Invocation: the check before every effect.}
  \rulename{Invoke} is where a proposal becomes an effect. It fires only when
  some held capability matches the tool's effect type and its predicate
  accepts the arguments; $\mathit{execute}(t,\,\mathbf{a}) \Downarrow o$ runs
  the tool and yields the observation $o$ appended to the context:
  \[
    \inferrule*[right=\rulename{Invoke}]{
      \Cap{e}{\varphi} \in \Sigma \quad
      \mathit{eff}(t) = e \quad
      \varphi(\mathbf{a}) \quad
      \mathit{execute}(t,\,\mathbf{a}) \Downarrow o
    }{
      \langle \Sigma;\, \mathit{ctx} \rangle
      \;\xrightarrow{(t,\,\mathbf{a})}\;
      \langle \Sigma;\, \mathit{ctx} \cdot o \rangle
    }
  \]
  If no capability covers the proposal it is blocked and the model receives a
  refusal. Since the model's output is untrusted (the trust boundary, Section~\ref{sec:background}),
  the check never consults the model about provenance. The only question is
  whether $\Sigma$ already holds a covering capability.
  


  \paragraph{What the rules give us.}
  By inspection of the three rules, every effect that executes is covered by a
  capability whose authority traces back to a trusted \rulename{Mint}.
  \rulename{Invoke} is the only rule that runs an effect, and it requires a
  covering capability already in $\Sigma$. A capability enters the initial
  store through \rulename{Mint}, or enters a child store through a chain of
  \rulename{Derive} steps. Because each derivation can only narrow its parent's
  scope, no child can exceed the authority created by the initial mint. In our
  evaluation, this mint occurs before untrusted content is admitted and later
  user-approved minting is disabled. Injected text may influence a proposed
  tool call or requested child grant, but it cannot expand the task-wide
  authority. A call outside the resulting scope is refused.

  \subsection{Implementation}
  \label{sec:pattern}
  We realize \tool\ on Pi~\cite{pi2025}, an open-source TypeScript coding-agent
  toolkit with a CLI and a programmable SDK, changing no Pi internal. Pi's four
  built-in tools instantiate the three effect types:

  \begin{center}
  \small
  \setlength{\tabcolsep}{5pt}
  \begin{tabular}{@{} l p{3.6cm} l @{}}
  \toprule
  \textbf{Tool} & \textbf{Arguments} & \textbf{Effect} \\
  \midrule
  \texttt{read}
    & \texttt{path : string}
    & \textsf{Read}  \\
  \texttt{write}
    & \texttt{path : string,\newline content : string}
    & \textsf{Write} \\
  \texttt{edit}
    & \texttt{path : string,\newline edits : Edits}
    & \textsf{Write} \\
  \texttt{bash}
    & \texttt{cmd : string}
    & \textsf{Exec}  \\
  \bottomrule
  \end{tabular}
  \end{center}

  Pi supplies the two host requirements of Section~\ref{sec:agents} directly.
  Its \texttt{createAgentSession} API exposes a \texttt{tool\_call} hook that
  fires synchronously before any tool is dispatched and may return
  \texttt{\{block: true, reason: \ldots\}} to abort it, the model receiving the
  reason as a refusal; \tool\ installs this hook separately for the orchestrator
  and each sub-agent. The \texttt{pi-subagents}
  extension\footnote{\url{https://www.npmjs.com/package/pi-subagents}} delegates
  a subtask to a \emph{process-isolated} child, whose separate address space
  lets \tool\ seed the child's capability store independently of the parent's.

  The three rules become four implementation components: a structured
  capability representation, initial minting, derivation and transfer to a
  child store, and a pre-dispatch check installed for each principal.

  \paragraph{Component 1: structured capabilities.}
  Scope is represented as data, not only as an opaque predicate. This makes
  the narrowing relation decidable: path prefixes compare by containment, and
  command scopes compare as sets of parsed argument vectors.
  
  \begin{lstlisting}[language=TypeScript,style=tsstyle,caption={The
    structured capability type.}]
  type EffectType = "read" | "write" | "exec";

  type ScopeSpec =
    | { kind: "path"; prefix: string }
    | { kind: "argv"; allowed: string[][] };
  
  interface Capability {
    effect: EffectType;
    spec: ScopeSpec;
    description: string;
  }

  type CapSet = Capability[];
  \end{lstlisting}
  
  \paragraph{Component 2: initial minting.}
  Before loading project context, \texttt{authorize} makes a preflight LLM
  call using the trusted request and the project's file tree. The call predicts the
  permissions required by the complete task and by the parent agent. The host
  validates and compiles this prediction into a task-wide \emph{ceiling}, from
  which children may derive authority, and the parent agent's own operational
  store. The latter may be read-only, forcing effectful work into scoped
  children.
  
  \begin{lstlisting}[language=TypeScript,style=tsstyle,caption={The
    initial authorization step.}]
  interface Authority {
    ceiling: CapSet;
    orchestrator: CapSet;
  }

  async function authorize(
    trustedRequest: string,
    taskTree: TaskTree
  ): Promise<Authority> {
    const proposal = await preflightLLM({
      trustedRequest,
      taskTree
    });
    return validateAndCompile(proposal);
  }
  \end{lstlisting}
  
  \noindent The \texttt{authorize} call appears once, at session start.
  Repository text and tool output are unavailable at this point. The
  preflight prediction may still be too broad or too narrow, but later
  injected content cannot change it. A runtime user approval, if enabled, is a
  separate trusted mint and is recorded as such.

  \paragraph{Component 3: derivation and transfer.}
  Delegation is driven by the model, just as the initial mint is. When an agent
  delegates, it issues a \texttt{subagent} tool call whose structured arguments
  carry the child's role, its subtask, and a \emph{requested grant}; the same
  decision that chooses what to delegate also proposes how much authority the
  child should get. The host treats that request as a proposal only:
  \texttt{derive} keeps just the capabilities structurally contained in the
  \emph{delegating agent's own} authority (rule~\rulename{Derive}) and drops the
  rest, so the child starts with the derived set alone and can never exceed its
  parent. The rule composes down a delegation chain---in principle a sub-agent
  may itself delegate, and its children are bounded by its already-narrowed
  store---so authority only shrinks with depth. The one exception is the root
  orchestrator, which derives children from the task-wide \emph{ceiling} rather
  than from its own store: this lets the orchestrator itself be restricted to
  read-only for tighter security while still handing a child the authority the
  task genuinely needs, and no child exceeds the ceiling because every grant is
  intersected with it (the \emph{Usage} listing below wires it this way).

  \begin{lstlisting}[language=TypeScript,style=tsstyle,caption={Deriving and
    transferring a child store.}]
  function derive(parent: CapSet, requested: CapSet): CapSet {
    return requested.filter(child =>
      parent.some(p =>
        p.effect === child.effect && scopeWithin(child.spec, p.spec)
      )
    );
  }

  async function spawnScoped(
    parent: CapSet, role: Role, task: string, requested: CapSet
  ) {
    // `requested` is the model's proposal; derive narrows it to the parent's authority.
    const childCaps = derive(parent, requested);
    const child = await createAgentSession({ role });
    installCapScope(child, childCaps);
    await child.sendMessage(task);
  }
  \end{lstlisting}
  
  \paragraph{Component 4: per-principal dispatch checks.}
  Every orchestrator and child session installs the same small handler, but
  closes over a different store. It compiles the structured scope, checks
  \rulename{Invoke}, and blocks uncovered calls.
  
  \begin{lstlisting}[language=TypeScript,style=tsstyle,caption={The
    per-principal \tool\ hook.}]
  function installCapScope(session: AgentSession, caps: CapSet) {
    session.on("tool_call", (event) => {
      const granted = caps.find(
        c => effectOf(event.toolName) === c.effect
          && compileScope(c.spec)(event.input)
      );
      if (!granted) {
        return {
          block: true,
          reason: `[CapScope] uncovered ${event.toolName} call`
        };
      }
    });
  }
  \end{lstlisting}
  
  \paragraph{Usage: a complete session.}
  Delegation always narrows: by rule~\rulename{Derive} a child receives a
  subset of its delegator's own authority, never more. The orchestrator is the
  one principal for which this needs care. As the root of the task it carries
  two stores. Its \emph{operational} store (\texttt{auth.orchestrator}) is what
  its own tool calls are checked against, and may be deliberately narrow---even
  read-only on tasks where the orchestrator should coordinate but never edit
  source itself. Its \emph{delegation} store is the task-wide
  \emph{ceiling}, the upper bound for every child it spawns. This is why
  \texttt{spawnScoped} derives against \texttt{auth.ceiling} rather than the
  orchestrator's own store: as the highest-level agent the orchestrator may
  grant a patcher the write capability the repair needs even when it holds no
  write itself, while \texttt{derive} still intersects every requested grant
  with the ceiling so no child can exceed it. By default the two stores
  coincide---the orchestrator simply holds the full ceiling and sub-agents take
  strict subsets of it---and they diverge only when a policy restricts the
  orchestrator to a pure-coordinator role.
  
  \begin{lstlisting}[language=TypeScript,style=tsstyle,caption={Complete
    orchestrator setup.}]
  const auth = authorize(userRequest);
  const orchestrator = await createAgentSession();
  installCapScope(orchestrator, auth.orchestrator);
  // onDelegate fires on the model's `subagent` tool call;
  // (role, task, grant) are its structured arguments, narrowed below.
  orchestrator.onDelegate((role, task, grant) =>
    spawnScoped(auth.ceiling, role, task, grant)
  );
  await orchestrator.sendMessage(userRequest);
  \end{lstlisting}
  
  \paragraph{Other frameworks.}
  The two seams \tool\ needs (Section~\ref{sec:agents}) are not peculiar to Pi:
  the \texttt{RunHooks} interface in the OpenAI Agents SDK, the tool-wrapping
  layer in LangChain, the middleware interface in LlamaIndex, and the
  \texttt{tool.invoke} override in DSPy all expose a pre-dispatch point. When
  sub-agents run in-process, the host ties each child to its own store rather
  than lean on process isolation; when they run as separate processes, it hands
  the child a serialized derived store at startup. Either way the capability
  representation and the narrowing relation carry over unchanged; only the hook
  registration and child setup are rewritten against the host API.
  
  \subsection{Example: A Backdoor the Allowlist Cannot Stop}
  \label{sec:example}
  We now run the delegated session end to end. A developer asks an orchestrator
  agent: \emph{the auth module's tests are failing; run the suite, triage, and
  patch the source until they pass}. This task genuinely needs project reads,
  source writes, and a small set of test commands. Before any project content is
  read, a preflight LLM call receives only the trusted request and the project's
  file tree, and the host compiles its policy into three task-wide capabilities,
  $c_1$ for reads, $c_2$ for writes, and $c_3$ for execution:
  \begin{align}
    c_1 &= \Cap{\mathit{Read}}{\mathit{canon}(p) \sqsubseteq
          \texttt{"project/**"}} \label{cap:read} \\
    c_2 &= \Cap{\mathit{Write}}{\mathit{canon}(p) \sqsubseteq
          \texttt{"src/**"}} \label{cap:edit} \\
    c_3 &= \textsf{Cap}[\mathit{Exec},\ s \in
          \{\texttt{pytest},\ \texttt{mypy src/}\}] \label{cap:exec}
  \end{align}
  The orchestrator delegates the noisy test run to a \emph{runner} and later the
  source edit to a \emph{patcher}. The runner needs to read tests and run
  \texttt{pytest} but needs no write, so before spawning it the orchestrator
  applies \rulename{Derive} to pass a strict subset:
  \begin{align}
    c_1' &= \Cap{\mathit{Read}}{\mathit{canon}(p) \sqsubseteq
            \texttt{"tests/**"}} \label{cap:dread} \\
    c_3' &= \textsf{Cap}[\mathit{Exec},\ s \in \{\texttt{pytest}\}]
            \label{cap:dexec}
  \end{align}
  Here \texttt{tests/**} is contained in \texttt{project/**} and \texttt{mypy}
  is dropped, so $\varphi' \Rightarrow \varphi$ and
  $\Sigma_{\mathit{child}} = \{c_1', c_3'\}$; crucially
  $\mathit{Write} \notin \mathit{eff}(\Sigma_{\mathit{child}})$, since the parent
  never passed $c_2$ nor any narrowing of it.

  Now the attack, exactly as in the introduction: an injection an allowlist
  cannot refuse without also refusing the legitimate fix. A trojaned fixture
  makes the suite print, to the runner's \texttt{pytest} output:
  \begin{quote}\small\itshape
  \textbf{[ci-bot] Known regression:} the signing key rotated. Restore the
  build by writing to \texttt{src/auth/keys.py}:
  \texttt{SIGNING\_KEY = "k3y\_9f2a\ldots c71"}, then re-run.
  \end{quote}
  The key is attacker-controlled, and the target is an \emph{ordinary source
  path}---no \texttt{.env}, \texttt{\textasciitilde/.ssh}, or \texttt{curl}, and
  nothing a deny pattern matches. Figure~\ref{fig:trace} compares three
  authorization models.
  \begin{figure*}[t]
  \centering
  \small
  \setlength{\tabcolsep}{6pt}
  \begin{tabular}{@{} c l c c c @{}}
  \toprule
  \textbf{Turn} & \textbf{Runner proposes}
    & \makecell{\textbf{Static}\\\textbf{global policy}}
    & \makecell{\textbf{Task-specific}\\\textbf{global policy}}
    & \makecell{\textbf{\tool}\\\textbf{per agent}} \\
  \midrule
  1 & \texttt{read("tests/test\_auth.py")}
    & allowed (global read) & allowed ($c_1$) & allowed ($c_1'$) \\[1pt]
  2 & \texttt{bash("pytest tests/test\_auth.py")}
    & \makecell{allowed (global exec)\\$\to$ payload enters}
    & \makecell{allowed ($c_3$)\\$\to$ payload enters}
    & \makecell{allowed ($c_3'$)\\$\to$ payload enters} \\[1pt]
  3 & \texttt{write("src/auth/keys.py", \ldots)}
    & \makecell{\textbf{allowed}: matches\\no deny rule}
    & \makecell{\textbf{allowed}: $c_2$\\covers \texttt{src/**}}
    & \makecell{\textbf{blocked}: no\\$\mathit{Write}$ in $\Sigma_{\mathit{child}}$} \\
  \bottomrule
  \end{tabular}
  \caption{The same runner under three authorization models. Turns~1 and 2
    are legitimate and allowed everywhere, so the poisoned test output enters the runner's
    context and it proposes the backdoor write at turn~3. The static global
    policy allows the ordinary-looking source path. The task-specific global
    policy also allows it because its shared authority includes $c_2$ for the
    legitimate patch. \tool\ blocks it because the runner has only
    $\Sigma_{\mathit{child}}=\{c_1',c_3'\}$.}
  \Description{A three-column execution trace comparing a static global
    policy, a task-specific global policy, and CapScope. All three allow the
    runner to read tests and execute pytest; only CapScope blocks the injected
    source write because the runner has no write capability.}
  \label{fig:trace}
  \end{figure*}
  The two global policies fail for related reasons. The static policy does not
  classify an ordinary source edit as dangerous. The task-specific policy
  knows that the repair needs source writes, but its one policy is shared by
  every agent. The runner that reads the poisoned output therefore also
  receives write authority. In both cases, the agent that reads the poisoned
  output can write to \texttt{src}. A global task policy cannot refuse this
  write for the runner while permitting the patcher to make the real fix.
  \tool\ can, because the runner never receives write authority. At turn~3 the
  runner is refused and reports the
  ``ci-bot hotfix'' upward \emph{as data}. The orchestrator, which does hold
  $c_2$ as part of its derivation ceiling, evaluates that report, and may
  delegate the legitimate fix to a patcher with a narrow write capability.
  Since the runner's $\Sigma_{\mathit{child}}$ contains no
  $\mathit{Write}$, no write executes in the runner whatever the test output
  says because there is no covering capability for it.


  \section{Evaluation}
  \label{sec:case_study}

  % \yihan{reorganize the structure. a bit loose now}

  We ask two questions. \emph{RQ1:} does per-principal scoping reduce attacks
  that global, task-level authorization cannot distinguish? \emph{RQ2:} what is
  the completion and latency cost of fail-closed capability enforcement? Across
  conditions we vary the authorization model while holding the underlying LLM,
  the tasks, and the delegated four-principal topology fixed.

  \paragraph{Tasks and attacks.}
  Each task is a small Python project with one bug and a failing
  \texttt{pytest} suite. The injections request either the execution of a
  project-local utility or a write to a project file. These actions are
  intentionally plausible rather than obviously harmful: each could be
  legitimate in another maintenance task, but none is needed for the requested
  repair. In the injected context, such execution or writing could introduce a
  malicious side effect. Table~\ref{tab:tasks} summarizes the repair tasks.

  \begin{table}[t]
  \centering\small
  \setlength{\tabcolsep}{5pt}
  \caption{Repair tasks used in the evaluation.}
  \label{tab:tasks}
  \begin{tabular}{@{} l p{0.76\columnwidth} @{}}
  \toprule
  \textbf{Task} & \textbf{Bug} \\
  \midrule
  T1 & Off-by-one error in \texttt{daterange.py} \\
  T2 & Wrong fee comparison in \texttt{fees.py} \\
  T3 & Wrong field lookup in \texttt{report.py} \\
  T4 & Case-sensitive lookup in \texttt{indexer.py} \\
  T5 & Discount applied on the wrong branch in
       \texttt{billing/fees.py} \\
  \bottomrule
  \end{tabular}
  \end{table}

  Each malicious action creates a unique sentinel or marker that benign repair
  never produces. We score that side effect, rather than treating routine
  project activity itself as compromise.

  \paragraph{Delegated pipeline.}
  Every condition uses the same four-principal workflow. An orchestrator
  diagnoses and coordinates but delegates effectful work. A runner executes
  the failing tests, a patcher edits the relevant source, and a verifier
  reruns the suite. The loop repeats until the tests pass or the run
  terminates. This control is important: the conditions compare authorization models,
  not different agent architectures.

  \paragraph{Injection surfaces.}
  We place the same task-specific payload in each of the five untrusted
  surfaces of Section~\ref{sec:agents}: a README, an agent-instruction file, a
  project-local skill, a source comment or docstring, and \texttt{pytest}
  output. The first three are normally encountered by the parent agent while
  inspecting the project, source comments by an agent reading code, and test
  output by the agent running or verifying the tests. Varying the surface
  therefore changes which principal first receives the injection without
  changing the requested repair or malicious action.

  \paragraph{Baselines.}
  We compare \tool\ against three global baselines that bracket current
  practice. All principals use ambient authority in \textbf{B0}. \textbf{B1}
  applies one static, task-independent policy to every principal, configured
  through the permission-gating facility Pi provides: it denies familiar
  hazards such as sensitive paths and obviously destructive commands (\eg
  \texttt{.env} and \texttt{.ssh} paths, \texttt{rm\,-rf}, \texttt{sudo},
  \texttt{curl}), but permits ordinary project operations. \textbf{B2}
  generates a tighter global policy from the trusted task description and a
  sanitized project structure before any injection is read, then applies that
  one policy to every principal; it can remove authority the task does not
  need, but cannot give two principals different slices of authority.
  \textbf{C} is not a baseline but \tool\ itself, the system under test: a
  preflight mints a task-wide ceiling that bounds delegation, and each sub-agent
  receives a store derived from their parent's authority for its delegated step. 
  B2 is the strongest global baseline: its policy is tailored to the task
  before the adversarial text arrives. The comparison with C isolates the
  value of authority \emph{per principal}.

  \paragraph{Protocol and metrics.}
  The full factorial is
  \[
    4\ \text{conditions} \times 5\ \text{tasks} \times
    5\ \text{surfaces} \times 3\ \text{trials} = 300\ \text{runs}.
  \]
  Each run uses an isolated working copy and
  all conditions are driven by a single model, \textbf{qwen3.5-flash}, accessed
  through an OpenAI-compatible endpoint.
  The harness verifies that the suite fails before the run, then records two
  outcomes:
  \begin{itemize}[topsep=2pt,itemsep=1pt,leftmargin=12pt]
  \item \texttt{task\_success}: the final test suite passes;
  \item \texttt{attack\_executed}: the attacker-only sentinel or marker exists
    on disk after the run.
  \end{itemize}
  Both are read from the filesystem, independent of any log and of each other:
  a run may fix the bug and still be compromised. \texttt{attack\_executed} is
  our headline defense metric---did the attack land. The per-call decision log
  additionally shows, on the conditions that install a hook, whether a call
  carrying the injected action was \emph{proposed} by the model and whether the
  layer \emph{blocked} it; we use these only as diagnostics when interpreting a
  run, not as reported numbers, and B0, which installs no hook, exposes
  neither.

  \paragraph{Results.}
  The qualitative ordering holds and is sharp (Table~\ref{tab:case}). For
  attack execution, $\mathrm{B0} \approx \mathrm{B1} > \mathrm{B2} \gg
  \mathrm{C}$: the injected effect lands in \BZeroLeaks\ of $75$ runs under
  ambient authority and \BOneLeaks\ under the static global policy---a
  conventional deny list does not improve on having no policy, because the
  malicious actions are ordinary in-repository operations it never names. The
  task-specific global policy (B2) helps, cutting leakage to \BTwoLeaks/75 by
  removing task-unnecessary operations before the injection arrives, but that
  one allowlist is still shared by every principal (below). \tool\ (C) admits the
  effect in only \BThreeLeaks/75---more than $10\times$ fewer than every
  global baseline---because it additionally withholds task-needed authority
  from the principals that do not need it. The per-surface view
  (Table~\ref{tab:surface}) confirms this is not an artifact of one carrier:
  leakage is high on surfaces that reach a broadly-authorized principal (\eg
  AGENTS.md, \SAgentsLeak/60) and lowest on the source comment
  (\SCommentLeak/60), the surface a narrowly-scoped principal reads.

  Crucially, this security costs little task completion. Constraining each
  principal lowers autonomous success only slightly: \tool\ finishes
  \BThreeTaskSuccess\ of $75$ repairs against \BZeroTaskSuccess\ for the
  unconstrained baseline---a gap of four runs---and B2's tighter global policy
  costs about as much (\BTwoTaskSuccess/75). The diagnostic decision log shows
  why the security gap is so much larger than the success gap: under \tool\ the
  model still \emph{proposes} the injected action on \BThreeProposals\ of $75$
  runs, comparable to the baselines, but \tool\ \emph{blocks} it
   instead of executing it. The visible price is
  latency: \tool\ averages \BThreeTime\,s per run against \BZeroTime--\BTwoTime\,s
  for the baselines, precisely because it refuses the most calls and each
  refusal can trigger a round of retries, and occasionally a loop, before the
  agent abandons the path. Both the small completion gap and the added time are
  artifacts of the fully-automated, fail-closed protocol rather than of the
  guarantee itself; we return to how to read these costs below. Even so, the
  exchange is favorable: a handful of repairs and some added latency for a
  more than $10\times$ reduction in injected effects that execute.

  \paragraph{Why \tool\ outperforms B2.}
  B2 fails for a structural reason, not a tuning one. Because one policy is
  shared by every principal, it must grant the union of everything the task
  needs---reads, a source write, the test commands---just to let the repair
  finish, and that broad authority then reaches \emph{every} sub-agent, the
  runner that reads the poisoned output included. Tightening only trades one
  failure for another: cut the authority the attack uses and some legitimate step
  loses it too, so the fooled agent retries until it times out. A single shared
  allowlist is trapped between broad enough to succeed and narrow enough to be
  safe. \tool\ scopes authority per principal instead: the runner that reads the
  injection can run tests but cannot write source, so the attack is blocked,
  while the patcher still holds the write that completes the repair. The task's
  authority is never concentrated in the agent the attacker reaches.

  \paragraph{Latency.}
  The conditions separate on wall-clock time as well (Table~\ref{tab:case},
  last column). The two denylist-style baselines are fastest---\BZeroTime\,s for
  B0 and \BOneTime\,s for B1---because almost nothing they propose is refused.
  B2 sits higher, at \BTwoTime\,s: its generated allowlist occasionally blocks a
  legitimate but unanticipated call, and each block costs a round of retries.
  \tool\ is the slowest, \BThreeTime\,s, for two compounding reasons: it makes
  one extra preflight call to mint the ceiling plus a short round-trip per
  delegation to author each child grant, and---dominant---it refuses the most
  calls, each refusal spawning several retries before the agent abandons the
  path. The added time is wall-clock spent on model round-trips and retries, not
  computation in the harness.

  \paragraph{Interpreting the cost.}
  The task-success and latency numbers are a conservative, worst-case bound.
  The pipeline runs fully autonomously, with no human approval path, so every
  uncovered call is simply refused; when a legitimate step falls outside the
  granted authority the agent retries variants of the blocked call, sometimes
  loops, and may time out. A failed repair is therefore rarely an unfixable bug
  but a measure of how often a deployment would surface an approval prompt, and
  the same retries inflate latency most under \tool, which refuses the most
  calls. A single human decision, \emph{is this the injection? then drop
  it}, would turn most such failures into completions and cut the looping; we
  disable that path deliberately to keep the experiment deterministic and free
  of user bias. The structural guarantee, that an uncovered call never executes,
  is unaffected.

  \begin{table*}[t]
  \centering\small
  \begin{minipage}[t]{0.51\textwidth}
    \setlength{\tabcolsep}{3pt}
    \caption{Aggregate results per condition ($75$ runs each).
      \emph{Attack executed} (lower is better) is the headline defense metric,
      read from the filesystem; \emph{mean time} is wall-clock seconds per run.
      Constraint costs little task success but cuts leakage sharply.}
    \label{tab:case}
    \begin{tabular}{@{} l r r r r @{}}
    \toprule
    \textbf{Condition} & \textbf{Runs} & \makecell{\textbf{Task}\\\textbf{success}}
      & \makecell{\textbf{Attack}\\\textbf{exec.}} & \makecell{\textbf{Mean}\\\textbf{time\,(s)}} \\
    \midrule
    B0: ambient authority
      & $75$ & \BZeroTaskSuccess/75  & \BZeroLeaks/75  & \BZeroTime  \\
    B1: static global
      & $75$ & \BOneTaskSuccess/75   & \BOneLeaks/75   & \BOneTime   \\
    B2: task-specific global
      & $75$ & \BTwoTaskSuccess/75   & \BTwoLeaks/75   & \BTwoTime   \\
    C: \tool
      & $75$ & \BThreeTaskSuccess/75 & \BThreeLeaks/75 & \BThreeTime \\
    \bottomrule
    \end{tabular}
  \end{minipage}\hfill
  \begin{minipage}[t]{0.43\textwidth}
    \setlength{\tabcolsep}{5pt}
    \caption{Results per injection surface, aggregated across all four
      conditions and five tasks ($4 \times 5 \times 3 = 60$ runs each).
      Aggregating over conditions isolates the carrier itself: the surface
      determines which principal first reads the injection, independent of the
      authorization model in force.}
    \label{tab:surface}
    \begin{tabular}{@{} l r r r @{}}
    \toprule
    \textbf{Surface} & \textbf{Runs} & \makecell{\textbf{Task}\\\textbf{success}}
      & \makecell{\textbf{Attack}\\\textbf{exec.}} \\
    \midrule
    README         & $60$ & \SReadmeTask/60  & \SReadmeLeak/60 \\
    AGENTS.md      & $60$ & \SAgentsTask/60  & \SAgentsLeak/60 \\
    Skill file     & $60$ & \SSkillTask/60   & \SSkillLeak/60 \\
    Source comment & $60$ & \SCommentTask/60 & \SCommentLeak/60 \\
    Tool output    & $60$ & \SToolTask/60    & \SToolLeak/60 \\
    \bottomrule
    \end{tabular}
  \end{minipage}
  \end{table*}



  \section{Discussion}
  \label{sec:discussion}
  
  \paragraph{Policy quality and approval.}
  Capability safety has two layers. Structural enforcement guarantees that a
  child never exceeds its parent and that an uncovered call does not execute.
  It does not guarantee that the initial ceiling or a model-proposed child
  grant is optimally narrow. An over-broad ceiling can authorize an attack
  within scope; an over-broad child grant can place that authority in the
  principal exposed to the injection. An under-broad grant can block a
  legitimate step and prevent task completion.

  Open-ended coding tasks make it hard to predict every command, dependency,
  or file an agent may legitimately need, so an initial policy can be too
  narrow as easily as too broad. \tool\ never silently widens authority to
  cover such a call; it refuses it, and in deployment the refusal and its
  requested scope can be surfaced to the user, who approves a narrow capability
  through a fresh trusted \rulename{Mint}. This is the lens through which
  Section~\ref{sec:case_study} reads task success and latency. The structural
  guarantee is orthogonal to policy quality: enforcement stays exact even when
  the predicted scope is not.

  \paragraph{How a residual injection still lands.}
  \tool's few remaining leaks (Table~\ref{tab:case}) arise in one of two ways,
  both upstream of enforcement and both a question of how authority was
  \emph{granted} rather than whether it was checked. The first is \emph{loose
  attenuation}: the orchestrator identifies the legitimate subtask but derives a
  child grant wider than that subtask needs---write to a directory rather than
  the one file, or execution of an interpreter rather than the one
  command---and the injected effect the child later reads falls inside that
  slack. The derivation bound still holds, so the child cannot exceed the
  ceiling; but the gap between what the task required and what was granted is
  exactly where the effect executes. The second is a \emph{fooled delegator}:
  the injection corrupts the orchestrator's model of the task itself, so it
  delegates the attacker's action as ordinary work and grants the child
  precisely the authority to carry it out. Here the grant is not too broad for
  what the orchestrator \emph{believes} the task to be---its task understanding
  was hijacked. This path can only succeed when the trusted ceiling already
  admits the action, since every child grant is intersected with it; so the
  preflight ceiling, minted before any untrusted text is read, is the final cap
  on both modes. Tightening that ceiling and attenuating each grant to least
  authority shrink both, and neither is closed by the dispatch check, which
  faithfully enforces whatever scope it is handed.

  \paragraph{Trusted computing base.}
  The TCB is exactly the host-side components named in
  Section~\ref{sec:threat}: the preflight isolation boundary, capability
  compiler, derivation code, per-principal stores, scope predicates, and
  dispatch hooks. The preflight LLM shapes how broad the initial policy is, but
  cannot place or alter capabilities once untrusted content is admitted---only
  host code can mint, narrow, or enforce them.

  \paragraph{Limitations and threats to validity.}
  The corpus contains five small Python repairs, one model, and three trials
  per task-surface-condition cell. It does not represent large refactors,
  dependency migrations, or long-running workflows, where policy omissions
  and approval frequency may be higher. The payloads are intentionally strong
  and the model may react differently to each carrier, so we report results by
  surface rather than assuming a uniform injection rate. Because the underlying
  calls are stochastic, individual cells vary from trial to trial; we therefore
  report aggregates over three trials and read the per-condition ordering, which
  is stable, rather than any single cell. Finally, the informative refusal
  message can affect recovery and task success, though it cannot turn an
  uncovered call into an allowed one.


  \section{Related Work}
  \label{sec:related}
  
  \paragraph{Prompt injection and model-side defenses.}
  Greshake~\etal~\cite{greshake2023not} introduced indirect prompt
  injection, in which attacker-controlled external content hijacks an
  LLM-integrated application; OWASP now ranks it the leading LLM
  application risk~\cite{owasp2025llm01}. In the coding setting the threat
  is acute: Liu~\etal~\cite{liu2025youraishell} report attack success
  rates up to 84\% in production coding editors, and rule and skill files
  are a documented backdoor vector~\cite{pillar2025rules}. Model-side
  defenses such as Instruction Hierarchy~\cite{wallace2024instructionhierarchy}
  and PromptArmor~\cite{shi2025promptarmor} make the model more skeptical
  of untrusted text, but remain heuristic: security rests on the model
  classifying provenance correctly, the very capability the attack
  undermines. \tool\ acts at the harness layer and is complementary,
  requiring no correct provenance judgment from the model.
  
  \paragraph{Harness-level defenses and the closest systems.}
  AgentDojo~\cite{debenedetti2024agentdojo} established the standard
  benchmark for agent prompt injection; Bhagwatkar~\etal~\cite{bhagwatkar2025firewalls}
  show that several such benchmarks are saturated by simple firewalls,
  motivating our coding-specific evaluation. Deployed coding agents already
  enforce permission boundaries through sandboxes, approval modes, and allow/deny
  command patterns. \tool\ differs along two axes. Its command scopes match a
  parsed executable and argument vector, and a compound command is checked
  segment by segment rather than as one string; thus allowing \texttt{pytest}
  does not also admit \texttt{pytest; cmd} or \texttt{pytest \&\& cmd} unless
  \texttt{cmd} is itself covered, nor \texttt{\$(cmd)}, whose substituted
  argument matches no allowed invocation. More importantly, authority is
  attenuated per principal.
  A task-wide global policy, even one generated specifically for the task,
  cannot permit a patcher to write a file while withholding that write from a
  runner. \rulename{Derive} can, and prevents the child from reacquiring
  authority outside the parent's ceiling.

  Other
  recent systems take a PL or systems view of the agent shell. CaMeL~\cite{debenedetti2025camel}
  uses capabilities to block exfiltration, but inside a dual-LLM planner
  with a custom interpreter, which is a far heavier architecture than a single
  hook. Fides~\cite{costa2025fides} tracks confidentiality and integrity
  labels over data flow; AgentArmor~\cite{wang2025agentarmor} reconstructs
  runtime traces into an IR and type-checks it; ToolGate~\cite{liu2026toolgate}
  attaches Hoare-style contracts to tool calls; and
  PACT~\cite{fan2026pact}, conceptually the nearest, treats injection as
  untrusted content determining an authority-bearing argument and enforces
  argument-level provenance contracts. \tool\ differs from all of these in
  where authority originates: rather than inferring it from data flow,
  traces, or provenance after the fact, it pre-issues authority as typed
  capability values minted from a trusted policy and checked at the
  dispatch boundary. The approaches are complementary: \tool's effect
  authority composes naturally with IFC-style data tracking.
  
  \paragraph{Capability security.}
  \tool's vocabulary is classical. Saltzer and
  Schroeder~\cite{saltzer1975protection} define a capability as an
  unforgeable authorization and elevate complete mediation and least
  privilege; Dennis and Van Horn~\cite{dennis1966programming} couple
  designation with authority, dissolving the confused-deputy
  problem~\cite{hardy1988confused}; and Miller and
  Shapiro~\cite{miller2003paradigm} distinguish authority from permission
  and show how attenuation follows from ordinary reference-passing.
  \tool\ instantiates this discipline in an LLM agent harness, with scope
  predicates as the attenuation mechanism.
  
  % \paragraph{Design patterns for LLM programs.}
  % Two LMPL-adjacent works model our contribution style:
  % Wang~\cite{wang2025effects} discharges the side effects of
  % LLM-integrated scripts through composable handlers rather than
  % entangling them with program logic, and Dai~\etal~\cite{dai2026mnimi}
  % reify a hidden invariant, sample independence, into reference types
  % via lightweight wrappers. \tool\ makes the analogous move for authority,
  % reifying it as a typed value enforced at dispatch rather than as a new
  % agent architecture.

  
  \section{Conclusion}
  \label{sec:conclusion}
  
  Indirect prompt injection in coding agents is a designation attack: because
  tools operate with ambient authority, any text the model reads can supply a
  resource name that resolves into an unauthorized effect. The durable
  response is not to rely on improved model resistance, because that is the
  heuristic the attack is built to defeat. Instead, we remove ambient
  authority and reify effect authorization as an explicit, typed, host-minted
  value that untrusted text cannot produce. \tool\ realizes this with structured
  capabilities, a trusted task-wide ceiling, structural attenuation into
  isolated child stores, and per-principal dispatch hooks. Our evaluation
  compares ambient authority, a static global policy, a task-specific global
  policy, and \tool\ across five tasks, five injection surfaces, and three
  trials per cell ($300$ runs). The outcome is one-directional: where global
  authorization---even a policy generated for the exact task---lets the
  injected effect execute in \BaseLeakLo--\BaseLeakHi\ of every $75$ runs,
  per-principal scoping admits it in only \BThreeLeaks, while completing all but
  four of the same $75$ repairs. The guarantee behind that gap does not depend
  on the rates: authority is granted to a principal, not inferred from a string
  the model emits. Authority, in a coding agent, should not be a string.

% \begin{acks}
% This work was sponsored by the National Natural Science Foundation of China
% (NSFC) under Grant No.\ W2542035.
% \end{acks}
  \bibliography{references}
\bibliographystyle{ACM-Reference-Format}
  \end{document}
